\documentclass[10pt,a4paper]{amsart}
\usepackage[margin=19mm]{geometry}
\usepackage{amsmath,amssymb,mathtools}
\usepackage[scr=rsfs,cal=euler]{mathalfa}
\usepackage{xfrac}
\usepackage[unicode,hidelinks,bookmarks=false]{hyperref}
\newcommand{\ie}{i.e.\ }
\newcommand{\cf}{cf.\ }
\newcommand{\resp}{respectively\ }
\newcommand{\Subsection}{\subsection}
\providecommand{\nobreakdash}{\nobreak\hbox{-}}
\setlength{\emergencystretch}{2em}
\multlinegap0pt
\usepackage{amssymb}
\usepackage{float}
\usepackage{mathtools}
\usepackage{xcolor}
\newtheorem{theorem}{Theorem}
\newtheorem{lemma}{Lemma}
\usepackage{cleveref}
\crefname{theorem}{Theorem}{Theorems}
\crefname{lemma}{Lemma}{Lemmas}
\title{Triangle-free subsets of the $r$-distance graph of the hypercube}
\author{Padmini Mukkamala, Ananthakrishnan Ravi}
\date{Source: arXiv:2506.18782v4 (CC BY 4.0)}
\begin{document}
\maketitle

\noindent\emph{Source and licence.} Triangle-free subsets of the $r$-distance graph of the hypercube. Version arXiv:2506.18782v4 (CC BY 4.0), distributed by arXiv under the Creative Commons Attribution 4.0 International licence, http://creativecommons.org/licenses/by/4.0/, as recorded in the arXiv licence metadata for this exact version and shown on the abstract page https://arxiv.org/abs/2506.18782v4. Full licence notice: \texttt{licenses/arxiv-2506.18782-CC-BY-4.0.txt}.\par\smallskip

\noindent\textit{Editorial scope.} Counted unit: the article's theorem thm:lower-main (source0.tex lines 126--177). The article's complete original proof of it is delivered with it (source0.tex lines 416--549, one \texttt{proof} environment of 134 lines, which the article places away from the statement). No mathematics is written, repaired or re-derived: the statement and the proof are the article's own lines, in the article's own notation, and no retained line is substituted or re-flowed.  Retained uncounted as the source-local context the proof needs: lines 284--290; lines 334--341; lines 388--394.  Nothing was omitted from any retained span, so the package carries no omission record.  No retained line contains a citation, so the wrapper carries no bibliography and no bibliography entry.  No retained line contains a figure environment, an image-inclusion command or drawing code, so no image is delivered and the package declares no asset dependency.  Editorial formatting only, in addition: an AMS \texttt{amsart} preamble that restates the article's own declarations and macros used by the retained lines, namely the environments \texttt{theorem}, \texttt{lemma} on separate counters, together with the standard packages those lines need.  The retained environments print in this document as Theorem 1, Lemma 1 in order of appearance, whereas the article prints the same items with the numbering of its own full document.  The complete native source of the article as distributed in the arXiv e-print for this exact version is bundled unchanged as \texttt{sources/180360/source0.tex}.
\begin{lemma}\label{lem:bch}
Let $d\ge 1$.
Then
\[
T(n,2d)\ \ge\ 2^{\,n-d\log_2(n+1)-d}.
\]
\end{lemma}

\begin{lemma}\label{lem:alteration}
Let $r$ be even and satisfy $r\le 2n/3$.
Then $H(n,r)$ has a triangle-free vertex set of size at least
\[
\frac{2\sqrt{2}}{3}\cdot
\frac{2^n}{\sqrt{\binom{n}{r}\binom{r}{r/2}\binom{n-r}{r/2}}}.
\]
\end{lemma}

\begin{lemma}\label{lem:twobit}
Let $r$ be even.
Then $H(n,r)$ has a triangle-free vertex set of size at least
\[
\binom{n}{r/2}-\binom{n-2}{r/2}.
\]
\end{lemma}

\begin{theorem}\label{thm:lower-main}
Let $r$ be even and satisfy $r\le 2n/3$.
Then $T(n,r)$ satisfies the following lower bounds.

\begin{enumerate}
\item Fixed distance $r=2d$.
For every fixed integer $d\ge 1$,
\[
T(n,2d)=\Omega_d\!\left(\frac{2^n}{n^d}\right).
\]

\item If $r=o(n^{1/3})$ with $r\to\infty$, then
\[
T(n,r)\ \ge\ 2^{\,n-\frac r2\log_2 n-O(r)}.
\]

\item If $r=o(n)$ with $r\to\infty$, then
\[
T(n,r)\ \ge\ 2^{\,n-\frac{3}{4}r\log_2(n/r)-O(r)}.
\]

We further divide the case $r=\alpha n$ (where $n$ ranges over integers for which $r=\alpha n$ is an even integer) into two intervals, one where $\alpha\in(0,\alpha_0)$ and the other where $\alpha\in[\alpha_0,2/3]$.
Numerically, $\alpha_0\approx 0.16241686$.

\item If $r=\alpha n$ with fixed $0<\alpha<\alpha_0$, then
\[
T(n,\alpha n)\ \ge\ 2^{(c(\alpha)-o(1))n},
\]
where
\[
c(\alpha)=
1-\frac12\left(
H_2(\alpha)+\alpha+(1-\alpha)H_2\!\left(\frac{\alpha}{2(1-\alpha)}\right)
\right).
\]

\item If $r=\alpha n$ with fixed $\alpha_0\le \alpha\le 2/3$, then
\[
T(n,\alpha n)\ \ge\ 2^{(H_2(\alpha/2)-o(1))n},
\]
where
\[
H_2(t)=-t\log_2 t-(1-t)\log_2(1-t)
\]
is the binary entropy function.
\end{enumerate}

Parts (1)--(2) follow from \Cref{lem:bch}.
Parts (3)--(4) are asymptotic corollaries of the alteration bound proved in \Cref{lem:alteration}.
Part (5) follows from the two-bit construction in \Cref{lem:twobit}.
Here, $\alpha_0\in(0,2/3)$ is the threshold at which the exponents of \Cref{lem:alteration} and \Cref{lem:twobit} agree.
\end{theorem}

\begin{proof}[Proof of \Cref{thm:lower-main}]
The five parts of \Cref{thm:lower-main} are proved by the three mechanisms established above: the minimum-distance code in \Cref{lem:bch}, the alteration bound in \Cref{lem:alteration}, and the two-bit construction in \Cref{lem:twobit}.
Hence $T(n,r)$ is at least the maximum of the relevant corresponding quantities.

\Cref{lem:bch} gives
\[
T(n,r)\ge 2^{\,n-\frac r2\lceil \log_2(n+1)\rceil}.
\]
For fixed $r=2d$, this implies
\[
T(n,2d)
\ge 2^{\,n-d\log_2(n+1)-d}
=
\Omega_d\!\left(\frac{2^n}{n^d}\right).
\]
This proves (1).

If $r=o(n^{1/3})$ and $r\to\infty$, then
\[
T(n,r)\ge 2^{\,n-\frac r2\lceil \log_2(n+1)\rceil}
       =2^{\,n-\frac r2\log_2 n-O(r)}.
\]
This proves (2).

We now analyze the alteration bound.
Using
\[
\binom{n}{r}\le \left(\frac{en}{r}\right)^r,
\qquad
\binom{r}{r/2}\le 2^r,
\qquad
\binom{n-r}{r/2}\le \left(\frac{2e(n-r)}{r}\right)^{r/2},
\]
\Cref{lem:alteration} gives, when $r=o(n)$ and $r\to\infty$,
\[
T(n,r)
\ge
2^{\,n-\frac{3}{4}r\log_2(n/r)-O(r)}.
\]
This proves (3).

When $r=\alpha n$ with fixed $0<\alpha<2/3$, we use
\[
\binom{n}{r}
\le
\frac{1}{\sqrt{2\pi n\alpha(1-\alpha)}}2^{nH_2(\alpha)}
\le
c\,2^{nH_2(\alpha)},
\]
where
\[
H_2(t)=-t\log_2 t-(1-t)\log_2(1-t)
\]
is the binary entropy function, and $c$ is a constant.
Then \Cref{lem:alteration} gives
\[
T(n,\alpha n)
\ge
2^{\left(
1-\frac12\left[
H_2(\alpha)+\alpha+(1-\alpha)H_2\!\left(\frac{\alpha}{2(1-\alpha)}\right)
\right]
-o(1)\right)n}.
\]
This proves (4).

Finally, \Cref{lem:twobit} gives
\[
T(n,r)
\ge
\binom{n}{r/2}\left(1-\frac{(n-r/2)^2}{n^2}\right).
\]
For $r=\alpha n$, we use
\[
\binom{n}{r/2}
\ge
\frac{2^{nH_2(\alpha/2)}}{n+1}.
\]
Therefore
\[
T(n,\alpha n)\ge 2^{(H_2(\alpha/2)-o(1))n}.
\]
This proves (5) whenever the two-bit exponent dominates.

To find the precise constant $\alpha_0$ which separates the two intervals where the above terms dominate, we define
\[
D(\alpha)
=
1-\frac12\left(
H_2(\alpha)+\alpha
+(1-\alpha)H_2\!\left(\frac{\alpha}{2(1-\alpha)}\right)
\right)
-H_2(\alpha/2).
\]
Putting $x=\alpha/2$, direct differentiation gives
\[
D'(\alpha)
=
\frac12\log_2\left(
\frac{x^5}{(1-3x)^3(1-x)^2}
\right).
\]
The function
\[
x\longmapsto
\frac{x^5}{(1-3x)^3(1-x)^2}
\]
is strictly increasing on $(0,1/3)$. Thus $D'$ changes sign at
most once. 
Moreover,
\[
\lim_{\alpha\to0^+}D(\alpha)=1>0,
\qquad
D(2/3)
=
1-\frac12\left(H_2(2/3)+\frac23\right)
-H_2(1/3)
<0.
\]
Since $D'$ changes sign at most once, from negative to positive,
the function $D$ first decreases and then increases. Its value at
the right endpoint is still negative, so it cannot cross zero on
the increasing part. It therefore has exactly one zero in
$(0,2/3)$.
We denote this zero by
$\alpha_0$. Numerically,
\[
\alpha_0\approx0.16241686.
\]
Consequently, the alteration exponent dominates for
$0<\alpha<\alpha_0$, whereas the two-bit exponent dominates for
$\alpha_0\le\alpha\le2/3$.

\end{proof}

\end{document}
